\chapter{Zadatak 3 --- diskretizacija konačnim razlikama}
\label{ch:z3-diskretizacija}

\section{Što se ovdje radi}

Nepoznata \emph{funkcija} $y(\cdot)$ zamjenjuje se konačnim vektorom njezinih
vrijednosti u čvorovima, a derivacije podijeljenim razlikama. Postupak je zadan
uputama, ali se ovdje izvodi do matričnog zapisa, koji upute ne daju.

\begin{predznanje}
\begin{itemize}[nosep,leftmargin=*]
  \item Kontinuirani problem \eqref{eq:z3-kont}--\eqref{eq:z3-kont-ogr}.
  \item Množenje matrice i vektora --- \S\ref{sec:mnozenje}.
  \item Oznaka $\coloneqq$ --- \S\ref{sec:znakovi-logike}.
\end{itemize}
\end{predznanje}

\section{Mreža}
\label{sec:z3-mreza}

\begin{korak}[Čvorovi i korak]
Interval $[0, L]$ dijeli se na $N$ jednakih podintervala duljine
\begin{equation}
\Delta x = \frac{L}{N},
\end{equation}
čime nastaje mreža od $N+1$ čvorova
\begin{equation}
x_i = i\,\Delta x, \quad i = 0, 1, \dots, N,
\qquad x_0 = 0, \quad x_N = L .
\end{equation}
Visine terena u čvorovima su $h_i = h(x_i)$, a nepoznanice su visine nivelete
\begin{equation}
y = [\,y_0,\ y_1,\ \dots,\ y_N\,]\trans \in \R^{N+1} .
\end{equation}
Broj \emph{čvorova} je $N+1$, broj \emph{podintervala} je $N$ --- razlika koja se
često promaši u dimenzijama matrica.
\end{korak}

\begin{zamka}[Zamka: $\Delta h$ nasuprot $\Delta x$ u tekstu zadatka]
Upute u istoj rečenici pišu ,,svaki duljine $\Delta h = L/N$``, a zatim u
formulama za razlike koriste $\frac{1}{\Delta x}$. To su iste veličine pod dvama
imenima --- vjerojatno omaška u tekstu. Ovdje se dosljedno piše $\Delta x$. Ako
te se na obrani pita, reci da si uočio nedosljednost u oznakama i uzeo
$\Delta h \equiv \Delta x = L/N$; to je jedino čitanje koje ima smisla, jer
$\Delta h$ inače ne bi bilo definirano nigdje drugdje.
\end{zamka}

\section{Podijeljene razlike}
\label{sec:z3-razlike}

Prema uputama zadatka:
\begin{formula}{Prva i druga podijeljena razlika}
\begin{align}
\nabla y_i &\coloneqq \frac{y_i - y_{i-1}}{\Delta x}, & i &= 1, \dots, N, \label{eq:z3-D1} \\
\nabla^2 y_i &\coloneqq \frac{\nabla y_i - \nabla y_{i-1}}{\Delta x}
             = \frac{y_i - 2y_{i-1} + y_{i-2}}{\Delta x^2}, & i &= 2, \dots, N. \label{eq:z3-D2}
\end{align}
\end{formula}

\begin{korak}[Odakle drugi izraz u \eqref{eq:z3-D2}]
Uvrštavanjem definicije prve razlike:
\[
\nabla^2 y_i
= \frac{1}{\Delta x}\left( \frac{y_i - y_{i-1}}{\Delta x} - \frac{y_{i-1} - y_{i-2}}{\Delta x} \right)
= \frac{y_i - y_{i-1} - y_{i-1} + y_{i-2}}{\Delta x^2}
= \frac{y_i - 2y_{i-1} + y_{i-2}}{\Delta x^2}.
\]
Koeficijenti $1, -2, 1$ karakteristični su za drugu razliku; zbroj im je nula,
što znači da druga razlika iščezava na svakoj \emph{linearnoj} funkciji --- kako
i treba, jer je druga derivacija pravca nula.
\end{korak}

\begin{zamka}[Zamka: koliko ograničenja stvarno ima]
\begin{itemize}[nosep,leftmargin=*]
  \item Prva razlika traži \emph{dva} uzastopna čvora, pa je definirana za
        $i = 1, \dots, N$ --- dakle $N$ vrijednosti, ne $N+1$.
  \item Druga razlika traži \emph{tri} uzastopna čvora, pa je definirana za
        $i = 2, \dots, N$ --- dakle $N-1$ vrijednosti.
\end{itemize}
Ako se to previdi, matrice $D_1$ i $D_2$ dobiju krive dimenzije i \texttt{linprog}
javi grešku o nepodudarnim veličinama. Ovo je najčešći uzrok pada koda u ovom
zadatku.
\end{zamka}

\begin{zamka}[Zamka: razlika unatrag, ne centralna]
Zadatak izričito propisuje razliku \emph{unatrag} --- $\nabla y_i$ koristi $y_i$
i $y_{i-1}$. Centralna razlika, $(y_{i+1} - y_{i-1})/(2\Delta x)$, bila bi
točnija (pogreška $O(\Delta x^2)$ umjesto $O(\Delta x)$), ali \emph{nije}
zadana. Ne mijenjaj propisanu shemu; ako je spomeneš, spomeni je kao alternativu
koju zadatak ne traži.
\end{zamka}

\section{Matrični zapis}
\label{sec:z3-matrice}

\begin{formula}{Matrice razlika}
$\nabla y = D_1 y$ i $\nabla^2 y = D_2 y$, gdje su
\begin{equation}
D_1 = \frac{1}{\Delta x}
\begin{bmatrix}
-1 & 1 & & & \\
& -1 & 1 & & \\
& & \ddots & \ddots & \\
& & & -1 & 1
\end{bmatrix} \in \R^{N \times (N+1)},
\end{equation}
\begin{equation}
D_2 = \frac{1}{\Delta x^2}
\begin{bmatrix}
1 & -2 & 1 & & \\
& 1 & -2 & 1 & \\
& & \ddots & \ddots & \ddots \\
& & & 1 & -2 & 1
\end{bmatrix} \in \R^{(N-1) \times (N+1)} .
\end{equation}
\end{formula}

\begin{korak}[Provjera dimenzija]
\begin{itemize}[nosep,leftmargin=*]
  \item $D_1$ ima $N$ redaka (po jedan za svaku prvu razliku) i $N+1$ stupaca
        (po jedan za svaki čvor). Umnožak $D_1 y \in \R^N$. \checkmark
  \item $D_2$ ima $N-1$ redaka i $N+1$ stupaca; $D_2 y \in \R^{N-1}$. \checkmark
  \item Svaki redak od $D_1$ ima točno dva elementa različita od nule, svaki
        redak od $D_2$ točno tri. Matrice su rijetke (\emph{sparse}) --- za
        $N = 100$ od $100 \times 101 = 10\,100$ elemenata samo $200$ nije nula.
\end{itemize}
\end{korak}

\begin{korak}[Zašto je $D_2$ podijeljena s $\Delta x^2$, a ne s $\Delta x$]
Jer je druga derivacija drugog reda: dvaput se dijeli korakom. Provjera
dimenzija: $y$ je u metrima, $\Delta x$ u metrima, pa je $D_1 y$
bezdimenzijski (m/m), kako i treba za nagib, a $D_2 y$ ima jedinicu
$\mathrm{m}^{-1}$, kako i treba za $b_2$. Ako se ovdje promaši potencija, brojevi
za $N = 20$ i $N = 100$ razlikovat će se za faktor $5$ umjesto da budu
usporedivi --- i to je najlakši način da otkriješ grešku.
\end{korak}

\begin{ploca}[Kako $D_1$ i $D_2$ zapisati na ploči]
Ne prepisuj cijele matrice. Napiši prvi redak i naznači strukturu:
\begin{enumerate}[nosep,leftmargin=*]
  \item ,,$D_1$ ima u svakom retku $-1$ i $1$ na susjednim mjestima, sve
        podijeljeno s $\Delta x$; dimenzija $N \times (N{+}1)$.``
  \item ,,$D_2$ ima $1, -2, 1$, podijeljeno s $\Delta x^2$; dimenzija
        $(N{-}1) \times (N{+}1)$.``
  \item Nacrtaj shematski pravokutnik s dijagonalnom vrpcom i upiši te brojeve u
        prvi redak. To je dovoljno --- ispitivač traži strukturu, ne prijepis.
\end{enumerate}
\end{ploca}
